\documentclass[12pt]{article}

%-------------------------
% PACKAGES
%-------------------------
\usepackage{amsmath} 
\usepackage{tikz}
\usepackage{ifthen}
\usepackage{fancyhdr}
\usepackage[margin=1in]{geometry}
\usepackage{xstring}
\usepackage{needspace} 
\usepackage{tabularx}   % For flexible table widths
\usepackage{booktabs}   % For professional-looking tables
\usepackage{xcolor}     % For coloring text (optional)
\usepackage{array}      % For defining custom column types

%-------------------------
% FBOX PARAMETERS
%-------------------------
% Reduce padding and border thickness for slimmer boxes
\setlength{\fboxsep}{3pt}    % Space between content and box
\setlength{\fboxrule}{0.5pt} % Thickness of the box border

%-------------------------
% CUSTOM COLUMN TYPES
%-------------------------
% Define a new column type 'L' for left-aligned, vertically centered content
\newcolumntype{L}[1]{>{\raggedright\arraybackslash}m{#1}}

%-------------------------
% HEADERS & FOOTERS
%-------------------------
\pagestyle{fancy}
\fancyhf{}

\setlength{\headheight}{28pt} 
\setlength{\headsep}{25pt}

\fancyhead[L]{\textbf{MATH 9999A\\Test 1}}
\fancyhead[C]{\textbf{SID:} \tikz \node[draw, rectangle, rounded corners, 
  minimum width=2in, minimum height=5mm] {};}
\fancyhead[R]{\textbf{VERSION: A\\October 15, 2022}}
\renewcommand{\headrulewidth}{2pt}
\fancyfoot[C]{\thepage}

%-------------------------
% QUESTION & OPTIONS MACROS
%-------------------------
% Initialize option commands as commands with one argument
\newcommand{\optionA}[1]{\def\optionAContent{#1}}
\newcommand{\optionB}[1]{\def\optionBContent{#1}}
\newcommand{\optionC}[1]{\def\optionCContent{#1}}
\newcommand{\optionD}[1]{\def\optionDContent{#1}}
\newcommand{\optionE}[1]{\def\optionEContent{#1}}

% Command to print the appropriate option based on the label
\newcommand{\printOption}[1]{%
  \ifthenelse{\equal{#1}{A}}{\optionAContent}{}%
  \ifthenelse{\equal{#1}{B}}{\optionBContent}{}%
  \ifthenelse{\equal{#1}{C}}{\optionCContent}{}%
  \ifthenelse{\equal{#1}{D}}{\optionDContent}{}%
  \ifthenelse{\equal{#1}{E}}{\optionEContent}{}%
}

%-------------------------
% LAYOUT MACROS
%-------------------------
% Normal Layout: Five options in a single row with fixed width
\newcommand{\fixedLabelsReorderedContentNormal}[1]{%
  % Extract each character from the reorder string
  \StrChar{#1}{1}[\charA]
  \StrChar{#1}{2}[\charB]
  \StrChar{#1}{3}[\charC]
  \StrChar{#1}{4}[\charD]
  \StrChar{#1}{5}[\charE]
  
  % Reduce space between columns for tight alignment
  \setlength{\tabcolsep}{2pt} % Default is 6pt
  
  % Adjust array rule width for thinner lines
  \setlength{\arrayrulewidth}{0.5pt}
  
  % Begin tabular environment with five left-aligned, vertically centered columns
  \begin{tabular}{|L{0.19\textwidth}|L{0.19\textwidth}|L{0.19\textwidth}|L{0.19\textwidth}|L{0.19\textwidth}|}
    \hline
    \textbf{A:}~\printOption{\charA} & 
    \textbf{B:}~\printOption{\charB} & 
    \textbf{C:}~\printOption{\charC} & 
    \textbf{D:}~\printOption{\charD} & 
    \textbf{E:}~\printOption{\charE} \\
    \hline
  \end{tabular}
}

% Enlarged Layout: Two columns per row to prevent empty rectangles
\newcommand{\fixedLabelsReorderedContentEnlarged}[1]{%
  % Extract each character from the reorder string
  \StrChar{#1}{1}[\charA]
  \StrChar{#1}{2}[\charB]
  \StrChar{#1}{3}[\charC]
  \StrChar{#1}{4}[\charD]
  \StrChar{#1}{5}[\charE]
  
  % Adjust tabular parameters for better alignment
  \setlength{\tabcolsep}{4pt} % Increased from 2pt to accommodate longer text
  \setlength{\arrayrulewidth}{0.5pt} % Thinner lines
  
  % Begin tabular environment with two left-aligned, vertically centered columns
  \begin{tabular}{|L{0.475\textwidth}|L{0.475\textwidth}|}
    \hline
    \textbf{A:}~\printOption{\charA} & 
    \textbf{B:}~\printOption{\charB} \\
    \hline
    \textbf{C:}~\printOption{\charC} & 
    \textbf{D:}~\printOption{\charD} \\
    \hline
    \multicolumn{2}{|L{0.95\textwidth}|}{\textbf{E:}~\printOption{\charE}} \\
    \hline
  \end{tabular}
}

% Single Long Layout: One option spans the full width, remaining options aligned vertically below
\newcommand{\fixedLabelsReorderedContentSingleLong}[1]{%
  % Extract each character from the reorder string
  \StrChar{#1}{1}[\charA]
  \StrChar{#1}{2}[\charB]
  \StrChar{#1}{3}[\charC]
  \StrChar{#1}{4}[\charD]
  \StrChar{#1}{5}[\charE]
  
  % Reduce space between columns for tight alignment
  \setlength{\tabcolsep}{2pt} % Default is 6pt
  
  % Adjust array rule width for thinner lines
  \setlength{\arrayrulewidth}{0.5pt}
  
  % Begin tabular environment with single column spanning full width
  \begin{tabular}{|L{0.95\textwidth}|}
    \hline
    \textbf{A:}~\printOption{\charA} \\
    \hline
    \textbf{B:}~\printOption{\charB} \\
    \hline
    \textbf{C:}~\printOption{\charC} \\
    \hline
    \textbf{D:}~\printOption{\charD} \\
    \hline
    \textbf{E:}~\printOption{\charE} \\
    \hline
  \end{tabular}
}

%-------------------------
% MULTIPLE CHOICE QUESTION MACROS
%-------------------------
% Macro for normal questions with five options in one row
\newcounter{question}
\newcommand{\makeMCQ}[2]{%
  \stepcounter{question}%
  \Needspace{15\baselineskip} % Ensures there's enough space for the question
  \begin{samepage} % Prevents page breaks within the question
    \noindent \textbf{\large Question \thequestion. #1} % Question text
    \vspace{0.5em} % Space between question and options
    \fixedLabelsReorderedContentNormal{#2} % Display options in normal layout
    \vspace{1em} % Space before the next question
  \end{samepage}
}

% Macro for enlarged questions with two columns per row
\newcommand{\makeMCQEnlarged}[2]{%
  \stepcounter{question}%
  \Needspace{25\baselineskip} % Ensures there's enough space for the question and two rows of options
  \begin{samepage} % Prevents page breaks within the question
    \noindent \textbf{\large Question \thequestion. #1} % Question text
    \vspace{0.5em} % Space between question and options
    \fixedLabelsReorderedContentEnlarged{#2} % Display options in enlarged layout
    \vspace{1em} % Space before the next question
  \end{samepage}
}

% Macro for questions with one super long option spanning the entire line
\newcommand{\makeMCQSingleLong}[2]{%
  \stepcounter{question}%
  \Needspace{35\baselineskip} % Ensures there's enough space for the question and all options
  \begin{samepage} % Prevents page breaks within the question
    \noindent \textbf{\large Question \thequestion. #1} % Question text
    \vspace{0.5em} % Space between question and options
    \fixedLabelsReorderedContentSingleLong{#2} % Display options with single long option
    \vspace{1em} % Space before the next question
  \end{samepage}
}

%-------------------------
% DOCUMENT BODY
%-------------------------
\begin{document}

%-----------------------
% Question 1 (Normal Layout)
%-----------------------
\optionA{Donald Trump}
\optionB{Joe Biden}
\optionC{Bill Clinton}
\optionD{George Bush}
\optionE{Barack Obama}
\makeMCQ{Who is the president of the USA?}{ABCDE}

%-----------------------
% Question 2 (Normal Layout)
%-----------------------
\optionA{1}
\optionB{2}
\optionC{3}
\optionD{4}
\optionE{5}
\makeMCQ{What is 2 + 2?}{CABDE}

%-----------------------
% Question 3 (Normal Layout)
%-----------------------
\optionA{Merge Sort}
\optionB{Quick Sort}
\optionC{Bubble Sort}
\optionD{Insertion Sort}
\optionE{Radix Sort}
\makeMCQ{Which sorting algorithm typically has \(O(n^2)\) average time complexity?}{CABDE}

%-----------------------
% Question 4 (Normal Layout)
%-----------------------
\optionA{\(x^3\)}
\optionB{\(1\)}
\optionC{\(2x\)}
\optionD{\(4x^2\)}
\optionE{\(\frac{1}{x}\)}
\makeMCQ{What is the derivative of \(x^2\)?}{ABCED}

%-----------------------
% Question 5 (Normal Layout)
%-----------------------
\optionA{\(\tfrac{1}{2} e^{2x}\)}
\optionB{\(2 e^{2x}\)}
\optionC{\(x e^{2x}\)}
\optionD{\(e^x\)}
\optionE{\(e^{2x}\)}
\makeMCQ{Compute \(\displaystyle \int e^{2x}\, dx.\)}{AECBD}

%-----------------------
% Question 6 (Normal Layout)
%-----------------------
\optionA{\(0\)}
\optionB{\(1\)}
\optionC{\(2\)}
\optionD{\(\infty\)}
\optionE{\(\text{Does not exist}\)}
\makeMCQ{Evaluate \(\displaystyle \lim_{x \to 0} \frac{\sin x}{x}.\)}{BACDE}

%-----------------------
% Question 7 (Enlarged Layout)
%-----------------------
\optionA{\((x - y)(x + y)\)}
\optionB{\((x - y)(x^2 + xy + y^2)\)}
\optionC{\((x + y)(x^2 - xy + y^2)\)}
\optionD{\((x^3 - y^3)(x - y)\)}
\optionE{\((x - y)^3\)}
\makeMCQEnlarged{Which is the correct factorization of \(x^3 - y^3\)?}{BCDAB}

%-----------------------
% Question 8 (Single Long Layout)
%-----------------------
\optionA{An exceptionally long option that exceeds half of the text width to test text wrapping within the designated rectangle.}
\optionB{Another lengthy option that requires the rectangle to be enlarged to accommodate the extended text without overflow.}
\optionC{A standard short option.}
\optionD{A moderately long option to ensure that the rectangle adjusts its size appropriately.}
\optionE{Yet another long option that tests the flexibility of the rectangle sizing in the LaTeX document.}
\makeMCQSingleLong{Evaluate the following statements regarding advanced calculus topics.}{ABCDE}

%-----------------------
% Additional Example Questions
%-----------------------

%-----------------------
% Question 9 (Normal Layout)
%-----------------------
\optionA{\(f(x) = x^2\)}
\optionB{\(f(x) = \sin x\)}
\optionC{\(f(x) = e^x\)}
\optionD{\(f(x) = \ln x\)}
\optionE{\(f(x) = \sqrt{x}\)}
\makeMCQ{Which function is its own derivative?}{ABCDE}

%-----------------------
% Question 10 (Enlarged Layout)
%-----------------------
\optionA{The integral of \(\frac{1}{x}\) is \(\ln |x| + C\).}
\optionB{The derivative of \(e^{x}\) is \(e^{x}\).}
\optionC{The limit of \(\frac{\sin x}{x}\) as \(x \to 0\) is 1.}
\optionD{The derivative of \(\ln x\) is \(\frac{1}{x}\).}
\optionE{The integral of \(x\) is \(\frac{1}{2}x^2 + C\).}
\makeMCQEnlarged{Which of the following statements are true about calculus functions?}{ABCDE}

%-----------------------
% Question 11 (Single Long Layout)
%-----------------------
\optionA{Given the function \(f(x) = x^3 - 3x + 2\), find all real roots and verify them by substitution.}
\optionB{Determine the convergence of the series \(\sum_{n=1}^{\infty} \frac{1}{n^2}\) using the p-test.}
\optionC{Compute the limit \(\lim_{x \to \infty} \frac{e^x}{x^n}\) for any positive integer \(n\).}
\optionD{Find the second derivative of \(f(x) = \sin(x)\) and interpret its geometric meaning.}
\optionE{Evaluate the definite integral \(\int_{0}^{\pi} \sin(x) \, dx\) and explain its significance.}
\makeMCQSingleLong{Answer the following advanced calculus problems.}{ABCDE}

\end{document}
